\subsection{QQQ return conventions and bounded resolution}
The fiscal-year discrepancy is not explained by citing a general data-vendor mechanism. \citet{schwarz2026} document that the CRSP SIZ-to-CIZ transition changes monthly returns, largely through payout timing. Legacy SIZ monthly returns reinvest distributions at month-end; CIZ compounds daily returns with reinvestment on ex-dates, as also described by \citet{french2026ciz}. For a period without a split, the ex-date convention is
\[
 1+R=\prod_{d=1}^D\frac{C_d+D_d}{C_{d-1}},
\]
where $D_d$ is the distribution per share on the ex-date. Corporate actions and units must be aligned before applying this identity. Our downloaded QQQ history has no split in the September 2008--September 2009 comparison window. These conventions do not imply that cash was available for reinvestment before its payment date.

A contemporaneous SEC exchange filing independently reports a September 30, 2008 QQQQ close of \$38.91 \citep[p.~4]{sec2008close}. The daily vendor reports the same opening endpoint and a September 30, 2009 close of \$42.25. Their price-only return is 8.5839\%, already above the issuer's 8.41\% reported market total return. Positive dividend entitlements cannot lower terminal wealth below this price-only bound on identical endpoints, whether held as cash or reinvested at positive prices. The 2010 issuer prospectus specifies consolidated closing prices and ex-date reinvestment \citep[pp.~90--91]{qqq2010prospectus}; substituting today's midpoint convention for the historical definition would be unwarranted.

\ResearchTable{QQQ fiscal-year return reconciliation}{tab:qqqfollowup}{qqq_followup.tex}{September 30, 2008 to September 30, 2009; returns in percent. Vendor dividend events are rounded to three decimals. The month-end alternative accrues dividends as receivables on ex-dates and reinvests at each month's final close. No payment dates are invented. NAV return is a separate economic basis.}

The adjusted-close factor $A_d/C_d$ is retained separately from quoted OHLC. Floating-point display precision is immaterial at the reported scale. If the frozen adjusted endpoints 33.54 and 36.62 were rounded to the nearest cent, the total-return interval is [9.1519\%, 9.2143\%]. This can accommodate the small frozen-versus-daily difference, but not the issuer gap. Rounded dividends and the ex-date versus month-end choice likewise cannot explain approximately 75 basis points. Exact issuer quote inputs remain unavailable. An endpoint-field inconsistency or report error is possible; neither is established. The issuer total is not used to back-solve or overwrite any price.

A fixed follow-up plan replaces the complete QQQ monthly return history with daily adjusted-close ratios while retaining every other input. Separately, the twelve-month horizon return is shifted by $\pm100$ basis points. For each sign, the corresponding cumulative gross factor is either distributed evenly in log terms across twelve months or placed entirely in each of the twelve individual months. All 26 perturbations rerun QQQ, R and M without changing parameters. These are counterfactual sensitivity inputs, not recovered observations, and do not bound unrelated proxy errors.
\ResearchTable{Full-history sensitivity to QQQ data uncertainty}{tab:qqqsensitivity}{qqq_sensitivity.tex}{January 2000--September 2026; percent. Intervals are minima and maxima across the 26 fixed perturbations, not confidence intervals. The alternate feed can share an upstream lineage with the original vendor, so it is not treated as a wholly independent market record.}

The original downside-versus-growth interpretation survives this finite audit. The precise cause of the issuer mismatch remains a documented limitation, rather than a reason to search indefinitely or to claim authenticated prices throughout the early proxy era.

\subsection{Observed daily bars and next-session execution sensitivity}
There is no universal academic convention that makes a reported opening bar a guaranteed fill. Monthly adjusted-close research with explicit proportional costs is common: \citet{devanathan2026}, for example, use daily adjusted closes and month-end rebalancing. That convention is a portfolio study, not opening-auction validation. We therefore retain the inherited monthly experiment and add a matched execution-clock sensitivity using actual daily bars \citep{yahoo2026daily}.

Each of QQQ, VTI, BND, GLD, BIL and SPY has all 3,708 expected common XNYS sessions from December 30, 2011 through September 30, 2026 in the downloaded panel. Dates, unique observations, positive prices, OHLC bounds and corporate-action records are checked. Actual ETF observations replace the four asset sleeves and BIL replaces modeled reserves. All 36-month estimation histories needed for R and M in the 2012-start experiment lie within observed ETF coverage. Earlier mutual-fund and indicative spot observations retain their monthly-only interpretation. No open is imputed for an unavailable instrument.

Let $C_{i,d}$ and $O_{i,d}$ be vendor historical close and open, and $A_{i,d}$ adjusted close. Define $F_{i,d}=A_{i,d}/C_{i,d}$ and $\widetilde O_{i,d}=O_{i,d}F_{i,d}$. The observed opening quote is preserved separately; $\widetilde O$ expresses it in total-return units. Targets, covariance, state and the no-trade decision are frozen at the previous completed month-end. Existing holdings bear $\widetilde O_{i,d}/A_{i,d-1}$ before the fee-consistent rebalance; new holdings then earn $A_{i,d}/\widetilde O_{i,d}$. Undeployed initial cash earns no overnight ETF return. Both branches use the same information set, prices, reinvestment convention and costs, isolating the decision-to-fill delay within this panel.

\ResearchTable{Daily-panel timing sensitivity at 10 basis points per dollar traded}{tab:dailyexecution}{daily_execution.tex}{CAGR, month-end maximum drawdown and daily-close maximum drawdown, in percent. The boundary-close branch remains an idealization; the next-open branch uses observed opening bars. BIL, not the inherited bill-rate proxy, supplies reserve returns. The second period ends December 2025. All methods retain their specifications; the ten-month trend rule is the inherited simpler control.}

\ResearchTable{Next-open scenarios at 25 basis points per dollar traded}{tab:dailycost25}{daily_cost25.tex}{Percent. M anticipates each cost in its optimization objective, so this is a cost-aware policy scenario, not a fee-only deduction from fixed trades. Higher stipulated costs can alter allocation enough to raise realized returns; no scenario is selected as a replacement policy.}

A second fixed cost of 25 basis points per dollar traded is reported above. Costs are stipulated aggregate allowances, not measured broker spreads or market impact. Four added checks independently reconcile the same-panel close branch to the original monthly engine, verify the no-cost buy-and-hold opening-price identity, perturb the next open without changing its already frozen target, and test self-financing fees. An observed bar is not proof of auction participation, available size or an executable limit order. The total-return-unit convention also does not model actual distribution payment dates, settlement, whole shares or taxes. Accordingly, this is an empirical execution sensitivity with a clear clock, not certification of a live investment strategy.
